\begin{pastexam}[label={px:ncku-110-cce-ds-01}]{成大電通所 110 學年度 第 1 題}
（是非改錯；30 分）

(\textbf{30 pts}) True or False. Identify all errors in the sentence if it is not correct. You can get the points only when you correct all errors.
\begin{enumerate}[label=\alph*.]
\item Floyd-Warshall algorithm applies the dynamic programming to find all pair shortest paths and can be applied to a graph with negative weights and negative cycles.
\item There are no back edges, no forward edges, and no cross edges after we apply the breadth-first search algorithm to an undirected graph.
\item A biconnected graph has only one biconnected component and we can use breath-first search algorithm to find biconnected components in a graph.
\item The access time of a hash table with open addressing is related to the load factor.
\item There is no need for parentheses for the prefix notation because there is no ambiguity.
\item There totally exists $n$-1 threads in a post-thread binary tree with $n$ nodes, where the left thread of a node points to its predecessor and the right thread of the node points to its successor according to the in-order.
\item DijKstra's algorithm is a greedy algorithm which can be used to find a single source shortest path. It allows a graph G(V, E) with negative weights. The time complexity is O(E log V) when we use an adjacency matrix to implement it.
\item The failure function of the Knuth-Morris-Pratt algorithm applied to the pattern ``abcabcacab'' is shown as follows:
\begin{center}
\begin{tabular}{cccccccccc}
a & b & c & a & b & c & a & c & a & b\\
0 & -1 & -1 & 0 & 1 & 2 & 0 & -1 & 0 & 1
\end{tabular}
\end{center}
\item A circular queue can be used to resolve the problem of necessary to move data if a queue is full, where the queue is implemented by an array with size $N$. But it can only store $N$-1 data in order to distinguish whether the queue is full or empty. We can pop a data from the location ``front'' and push a data to the location ``rear+1'', where ``front'' and ``rear'' are pointers to the queue.
\item If f(n)=$\sum_{i=1}^{n}\log i$ , then f(n)= $\theta$(log n ).
\end{enumerate}
\end{pastexam}

\begin{pastsoln}
答案：a.~F\quad b.~F\quad c.~F\quad d.~T\quad e.~T\quad f.~F\quad g.~F\quad h.~F\quad i.~T\quad j.~F。

十句分屬不同章：a、g 是資結書第~25~章，b 第~21~章，c 第~23~章，d 第~12~章，e、i 第~6~章，f 第~7~章，h 第~27~章，
j 第~1~章；放在第~25~章。 錯的句子先寫改正後的句子，再說明。

\textbf{a.~F。} 改成：「\ldots can be applied to a graph with negative weights \textbf{but no} negative cycles.」 Floyd--Warshall
是動態規劃沒錯（資結書第~25.6~節），也允許負權重的邊。 但有負環時，繞著負環多走一圈路徑就更短，最短路徑沒有定義；
演算法跑完會出現 $d[i][i] < 0$，只能用來\textbf{偵測}負環，算出來的距離沒有意義。

\textbf{b.~F。} 改成：「There are no back edges and no forward edges, \textbf{but there may be cross edges} \ldots」。 BFS 中每條邊
$\{u, v\}$ 都有 $|d[u] - d[v]| \le 1$（資結書第~21.3~節）。 若一條非樹邊連到祖先或子孫，例如 $v$ 是 $u$ 的子孫，則
$d[v] \ge d[u] + 1$，配合上式只能 $d[v] = d[u] + 1$，也就是 $u$ 是 $v$ 的父節點，這條邊就是樹邊，矛盾。 所以非樹邊兩
端互不為祖先子孫，都是交叉邊（同一層或相鄰兩層之間）。 交叉邊確實會出現：三角形 $1, 2, 3$ 從 $1$ 做 BFS，
$\{1,2\}$、$\{1,3\}$ 是樹邊，$\{2,3\}$ 是交叉邊。 （對無向圖做 \textbf{DFS} 才是只有樹邊和後向邊，資結書定理~21.8。）

\textbf{c.~F。} 改成：「\ldots we can use \textbf{depth}-first search algorithm to find biconnected components \ldots」（原文的 breath
是 breadth 的筆誤）。 前半句對：雙連通圖沒有割點，整張圖就是一個雙連通元件。 找雙連通元件要用 DFS 樹算
$\mathit{low}$ 值（資結書第~23.5、23.6~節）：它依賴 DFS 的「非樹邊只連到祖先」（後向邊），$\mathit{low}[v]$ 才能表示
「$v$ 的子樹能爬回多高的祖先」；BFS 樹的非樹邊是交叉邊（見 b），沒有這個性質。

\textbf{d.~T。} 開放定址的探測次數由載入因子 $\alpha = n/m$ 決定：在均勻雜湊的假設下，失敗搜尋的期望探測次數至多
$1/(1 - \alpha)$，成功搜尋至多 $\frac{1}{\alpha}\ln\frac{1}{1-\alpha}$（資結書定理~12.5、12.7）。 $\alpha$ 越接近 $1$ 越慢。

\textbf{e.~T。} 每個運算子的運算元個數固定時，前序式（波蘭式）由左往右讀，每個運算子後面接著它的兩個運算元，切法唯一，
不需要括號也沒有優先順序的問題（資結書第~6.3.5~節）。 後序式也一樣。

\textbf{f.~F。} 改成：「There totally exist $n$\textbf{+}1 threads in an \textbf{in}-threaded binary tree with $n$ nodes \ldots」。 兩個錯：
\begin{itemize}
\item 引線指向「中序」的前驅和後繼，這是\textbf{中序}引線二元樹（in-threaded），不是 post-thread。
\item $n$ 個節點的二元樹有 $2n$ 個指標欄，只有 $n - 1$ 個指向子節點（每個非根節點被指一次），所以有 $n + 1$ 個空指標，
全部改成引線（資結書第~7.5~節）。 常見課本（Horowitz）用一個表頭節點，中序第一個節點的左引線和最後一個節點的右引線都
指向表頭，引線共 $n + 1$ 條。
\end{itemize}
（若不用表頭節點，那兩條引線沒有東西可指，留成空指標，真正指到節點的引線才是 $n - 1$ 條；資結書就是這樣寫的。 這題
要改的主要是 post-thread；寫 $n + 1$ 時最好註明是有表頭節點的版本。）

\textbf{g.~F。} 改成：「\ldots It \textbf{does not allow} a graph G(V, E) with negative weights. The time complexity is O(E log V) when we use an
\textbf{adjacency list (with a binary heap)} to implement it.」（或保留相鄰矩陣，把時間改成 $O(V^2)$。） Dijkstra 是貪心法、求單一
起點最短路徑，這兩點對。 它每次確定距離最小的頂點後就不再更新，靠的是「邊權重非負，後面不可能再出現更短的路」，有負邊
時會出錯（資結書第~25.2.5~節）。 用相鄰矩陣時，找每個頂點的鄰居要掃一整列，總共 $\Theta(V^2)$，再用陣列挑最小，總時間
$O(V^2)$；$O(E \log V)$ 是相鄰串列配二元堆積的時間（資結書第~25.2~節）。

\textbf{h.~F。} 原題的失敗函數用的是 Horowitz 的寫法：$f(j)$ 是 $p_0 \cdots p_j$ 最長的「既是前綴又是後綴的真子字串」的
最後一個位置，沒有時是 $-1$；它等於資結書的 $\pi[j+1] - 1$（資結書第~27.3~節）。 先算 $\pi$：
\begin{center}
\small
\begin{tabular}{l*{10}{c}}
\toprule
$j$ & 0 & 1 & 2 & 3 & 4 & 5 & 6 & 7 & 8 & 9 \\
$p_j$ & a & b & c & a & b & c & a & c & a & b \\
\midrule
最長邊界的長度 $\pi[j+1]$ & 0 & 0 & 0 & 1 & 2 & 3 & 4 & 0 & 1 & 2 \\
正確的 $f(j)$ & $-1$ & $-1$ & $-1$ & 0 & 1 & 2 & \textbf{3} & $-1$ & 0 & 1 \\
題目給的 & \textbf{0} & $-1$ & $-1$ & 0 & 1 & 2 & \textbf{0} & $-1$ & 0 & 1 \\
\bottomrule
\end{tabular}
\end{center}
有兩格錯：$f(0)$ 應為 $-1$（單一字元沒有真的邊界）；$f(6)$ 應為 $3$（$p_0 \cdots p_6 = \mathtt{abcabca}$，前綴 $\mathtt{abca}$ 也是
後綴，長度 $4$）。 另外看 $j = 7$：$\mathtt{abcabcac}$ 結尾是 c，從 $f(6) = 3$ 退回 $f(3) = 0$、再退到 $f(0) = -1$ 都接不上，
所以 $f(7) = -1$，表上這格是對的。

\textbf{i.~T。} 依題意，$\mathit{front}$ 指向前端元素、$\mathit{rear}$ 指向尾端元素：取出時拿 $\mathit{front}$ 那格，加入時放在
$(\mathit{rear} + 1) \bmod N$。 用陣列做的一般佇列，$\mathit{rear}$ 到了陣列尾端就必須把資料往前搬，環狀佇列讓索引繞回
開頭，免去搬移；只靠 $\mathit{front}$、$\mathit{rear}$ 時，空和滿無法區分，所以留一格不用，最多放 $N - 1$ 個（資結書第~6.2.2~節）。
三件事都對。 要注意兩點：一、「rear+1」要取 $\bmod N$；二、這和 $\mathit{front}$、$\mathit{rear}$ 的定義有關，若照 Horowitz
的寫法（$\mathit{front}$ 指向前端元素的\textbf{前一格}），取出的位置是 $\mathit{front} + 1$，這句的前半要跟著改。 而「只能放
$N - 1$ 個」是指只用兩個索引的情形，多記元素個數或一個旗標就能放滿 $N$ 個。

\textbf{j.~F。} 改成：「\ldots then f(n) = $\theta$(\textbf{n} log n).」 上界：每一項 $\log i \le \log n$，$f(n) \le n \log n$。 下界：只看後一半，
\[
f(n) \ge \sum_{i = \lceil n/2 \rceil}^{n} \log i \ge \frac{n}{2} \log \frac{n}{2} = \Omega(n \log n).
\]
所以 $f(n) = \log n! = \Theta(n \log n)$（資結書第~1.6~節；也可以用 Stirling 公式），不是 $\Theta(\log n)$。
\end{pastsoln}
